\section{Introduction}\label{Chapter-1-Introduction}
\noindent The rapid growth of large language models (LLMs) has transformed how text is generated. While these systems increase productivity across many sectors, they also present significant challenges to information security, academic integrity, and digital trust. Automated systems can now generate human-like text at scale, making it easy to create deceptive reviews, fake news, or spam that bypasses traditional security check-points. Detecting machine-generated text has therefore become a critical area of study in modern natural language processing (NLP).

\section{The Generalization Gap and Closed-Set Limits}\label{Chapter-1-Generalization-Gap}
\noindent Traditional AI-generated text detectors are built using supervised machine learning models. While they achieve high accuracy (often exceeding $95\%$) on held-out test data from their training distribution, they suffer from a severe performance drop when tested on unseen topics (domain shift) or text generated by different models (generator shift). Supervised models assume a closed-world setup where every input must belong to the classes seen during training, rendering them fragile under real-world, in-the-wild shifts.

When an AI text detector is deployed in the wild, it will inevitably encounter text generated by previously unseen LLM architectures. A closed-set generator attributor (e.g. classifying text as GPT vs. LLaMA) will be forced to select one of the known classes, even if the text was generated by a completely unseen model family (e.g., OPT). An explicit open-set generator rejection layer addresses this by introducing an \texttt{UNKNOWN} class. By modeling the boundaries of known generator distributions, the system can reject anomalous, out-of-distribution text rather than making incorrect attributions.

\section{Motivation}\label{Chapter-1-Motivation}
\noindent Our motivation originates from the need to improve the reliability and generalizability of machine-generated text detectors in real-world scenarios:
\begin{enumerate} 
\item \textbf{Generalization Gap in AI Text Detection:} Most supervised AI detectors perform exceptionally well in "closed-lab" environments where training and test distributions match, but fail under distribution shifts.
\item \textbf{Closed-Set Attributor Limitations:} Existing supervised attribution systems operate under a closed-world assumption, forcing every machine-generated text into a known class (e.g., classifying a text as GPT or LLaMA even if it was generated by OPT).
\item \textbf{Need for Rejection Capability:} A robust detector must have the ability to flag text generated by a previously unseen LLM as \texttt{UNKNOWN} rather than making an incorrect assignment.
\end{enumerate}

\section{Research Objectives}\label{Chapter-1-Objectives}
\noindent The objectives of this project are:
\begin{itemize}
    \item To reproduce a manageable, multi-domain text detection benchmark based on the ACL 2024 MAGE paper.
    \item To implement multiple detection paradigms (lexical, rank-based, zero-shot, and contextual supervised).
    \item To build an explicit open-set generator rejection layer that models known generator embedding distributions.
    \item To evaluate the OOD rejection performance under representation (pooling), distance metric, and cross-generator generalization studies.
\end{itemize}